%----------
%	LÍNEAS FUTURAS
%----------

Los resultados de este trabajo abren varias vías de continuación, ordenadas de la más
inmediata a la más ambiciosa.

La primera se deriva directamente del principal hallazgo. Si buena parte del desacuerdo
mide distancia de criterio y no error, tiene poco sentido seguir tratando cada variable
como si tuviera una única interpretación correcta. Resultaría prometedor explorar
enfoques que no asuman una etiqueta de oro única por categoría, como el aprendizaje con
desacuerdo y las aproximaciones multietiqueta, que permitirían al sistema representar la
ambigüedad de las variables más disputadas en lugar de forzar un veredicto binario. En
la misma línea, incorporar una doble codificación de una parte del corpus, ausente en
este estudio (Sección~\ref{subsec:limitacion_fiabilidad}), permitiría medir la
fiabilidad intercodificadora humana y situar con precisión el techo real de la tarea,
hoy solo acotado de forma indirecta.

Una segunda vía es refinar el aprovechamiento del contexto experto por parte de los
modelos. Dado que la ganancia de las \textit{Agent Skills} depende de la capacidad del
modelo para explotar la carga bajo demanda, conviene estudiar estrategias de
\textit{few-shot} con ejemplos verificados, así como el ajuste fino (\textit{fine-tuning})
de modelos abiertos como \texttt{gemma4:e4b} sobre el propio corpus anotado. El modelo
local, competitivo en acuerdo y sin coste de tarifas
(Sección~\ref{subsec:iris_local}), es el candidato natural para esta especialización,
con la ventaja añadida de preservar la privacidad del material.

Una tercera vía atañe a la explicabilidad y a la validación del plano cualitativo. El
sistema ya produce, para cada predicción, la explicación y las evidencias literales que
la sustentan, pero ese marcado no ha podido evaluarse por falta de un marcado experto de
referencia. Construir ese conjunto de referencia, en el que las analistas señalen sobre
el texto los fragmentos que motivan cada variable, permitiría medir la calidad de las
evidencias, no solo de la etiqueta, y cerrar así la evaluación de las dos dimensiones que
un análisis de contenido experto necesita. Sobre esa misma base trazable cabe construir
un flujo con humano en el bucle, en el que la codificadora acepte, rechace o corrija cada
veredicto y esa retroalimentación sirva para reajustar los criterios y los umbrales por
variable, de forma coherente con el papel de asistente, y no de decisor, que se ha
argumentado para el sistema.

Por último, en el plano de la aplicación, el objetivo a medio plazo es la integración
del sistema en una herramienta de postedición dentro del proyecto IRIS, que ordene el
corpus, resalte los pasajes candidatos y facilite el trabajo de las analistas. Este
trabajo se enmarca en la colaboración del proyecto IRIS con RTVE, cuyo Observatorio de
Igualdad aporta parte de las guías de lenguaje que fundamentan el libro de códigos, de
modo que un escenario de aplicación natural es el apoyo a los flujos de revisión
editorial de una corporación de radiotelevisión pública, donde el seguimiento de la
representación de las mujeres forma parte de su mandato de servicio público. Esa
herramienta permitiría además extender el análisis a más variables del libro de códigos
y a otros medios y periodos, y ofrecería a instituciones y equipos de investigación un
instrumento reproducible para el seguimiento sistemático de la representación de las
mujeres en los medios. La combinación de un despliegue local respetuoso con la privacidad
y una salida auditable por personas expertas es, a nuestro juicio, la vía más realista
para llevar este trabajo del entorno experimental al uso real.
